% !TEX program = lualatex
% Arquivo: Revisao.tex

\documentclass[a4paper,12pt]{article}

% ===== Idioma =====
\usepackage[portuguese]{babel}
\addto\captionsportuguese{%
  \renewcommand{\contentsname}{Sumário}%
}

% ===== Fonte (LuaLaTeX) =====
\usepackage{fontspec}
\setmainfont{Latin Modern Roman}

% ===== Layout =====
\usepackage[top=3cm,bottom=2cm,left=2.5cm,right=1.8cm]{geometry}
\usepackage{indentfirst}
\usepackage{setspace}
\onehalfspacing
\setlength{\parindent}{1.25cm}
\setlength{\parskip}{0pt}

% ===== Tabelas / listas =====
\usepackage{booktabs}
\usepackage{enumitem}
\setlist[itemize]{noitemsep,topsep=3pt}

% ===== Figuras gerais =====
\usepackage{graphicx}
\usepackage{float}
\usepackage{pdflscape}
\graphicspath{{./}{./imagens/}{./figuras/}}

% ===== Links =====
\usepackage[hidelinks]{hyperref}
\usepackage{xurl}
\Urlmuskip=0mu plus 1mu

% ===== TikZ (opcional; não quebra se o styles não existir) =====
\usepackage[dvipsnames]{xcolor}
\usepackage{tikz}
\usetikzlibrary{shapes.geometric, arrows.meta, positioning}
\IfFileExists{flu1_styles.tex}{\input{flu1_styles}}{}

% ===== Citações ABNT (autor-data) + BibTeX =====
\usepackage[alf]{abntex2cite}
Uso no main:
  \input{caminho/tab_comparativo_IMO_ANP_landscape}

Requer no preâmbulo do main (sugestão mínima):
  \usepackage{pdflscape}   % ou lscape
  \usepackage{longtable}
  \usepackage{tabularx}
  \usepackage{booktabs}
  \usepackage{ragged2e}
  \usepackage{array}
  

\begin{landscape}

\setlength{\LTpre}{6pt}
\setlength{\LTpost}{6pt}

\renewcommand{\arraystretch}{1.15}

\footnotesize
\begin{longtable}{@{}
  >{\RaggedRight\arraybackslash}p{0.14\linewidth}
  >{\RaggedRight\arraybackslash}p{0.41\linewidth}
  >{\RaggedRight\arraybackslash}p{0.41\linewidth}
@{}}

\caption{Comparação entre IMO Res. A.672(16) e Resolução ANP nº 817/2020 (Anexo I) — afinidades e diferenças principais}
\label{tab:comparativo_imo_anp}\\

\toprule
\textbf{Tema} & \textbf{Pontos de maior afinidade} & \textbf{Maiores diferenças} \\
\midrule
\endfirsthead

\multicolumn{3}{@{}l@{}}{\textit{(continuação)} — Tabela \ref{tab:comparativo_imo_anp}}\\
\toprule
\textbf{Tema} & \textbf{Pontos de maior afinidade} & \textbf{Maiores diferenças} \\
\midrule
\endhead

\midrule
\multicolumn{3}{r@{}}{\textit{continua na próxima página}}\\
\endfoot

\bottomrule
\endlastfoot

Natureza e alcance &
Ambas orientam a decisão sobre remoção/permanência por critérios de segurança e proteção ambiental, buscando evitar riscos a terceiros e ao meio marinho. &
A IMO opera como diretriz internacional (recomendatória) para governos; a ANP estrutura um procedimento regulatório nacional com etapas, documentos, aprovações e responsabilidades formais. \\

Regra geral &
As duas partem do princípio de que instalações abandonadas/em desuso devem ser removidas para evitar obstáculos e riscos. &
A IMO condiciona e enquadra exceções com lógica de reporte/aviso no âmbito marítimo; a ANP explicita também a vedação ao alijamento e amarra o tema ao ciclo regulatório do descomissionamento (PDI/EJD e afins). \\

Exceções (remoção parcial / permanência \textit{in situ}) &
Admitidas quando tecnicamente justificáveis e quando a solução não gera interferência injustificável na navegação e/ou outros usos do mar. &
A IMO explicita motivos típicos (novo uso, inviabilidade técnica, custo extremo, risco inaceitável); a ANP exige comparação de alternativas com critérios mínimos e aprovação pelas autoridades competentes no contexto brasileiro. \\

Critérios de decisão (\textit{case-by-case}) &
Abordagem multicritério (segurança, ambiente, viabilidade). &
A IMO dá ênfase forte a navegação, deterioração, deslocamento e efeitos ambientais; a ANP define explicitamente um conjunto mínimo (técnico, ambiental, social, segurança e econômico) e afirma que nenhum critério isolado é decisivo. \\

Navegação: coluna d'água e sinalização &
Convergência no requisito prático de manter coluna d'água desobstruída (quando aplicável) e tratar remanescentes como risco a gerenciar. &
A IMO fixa o requisito e o tratamento como padrão internacional; a ANP detalha gatilhos e obrigações de sinalização/gestão no âmbito marítimo (Marinha/DPC/Capitanias), conforme o caso. \\

Avisos e divulgação a navegantes &
Ambas pressupõem comunicação e gestão do risco para terceiros. &
A IMO detalha aviso prévio e atualização em cartas/publicações náuticas; a ANP trata de comunicações formais (ANP, órgão ambiental e autoridade marítima) sobretudo quando há alteração do PDI aprovado. \\

Monitoramento pós-descomissionamento &
Se algo permanece, deve haver monitoramento do local e do risco. &
A IMO descreve monitoramento sob responsabilidade do Estado costeiro/arranjo definido; a ANP exige plano pós-descomissionamento baseado em risco e prevê relatório consolidado no encerramento. \\

Responsabilidade e capacidade financeira &
As duas buscam assegurar que exista parte responsável por passivos futuros e gestão/monitoramento. &
A IMO exige título legal inequívoco e capacidade financeira quando não há remoção total; a ANP explicita responsabilidade do contratado/operador e prevê execução de garantias em caso de descumprimento. \\

Padrões quantitativos e “gatilhos” &
Uso de requisitos objetivos para reduzir risco à navegação e facilitar decisão. &
A IMO trabalha com limiares por profundidade e massa e inclui exigência de projetar novas instalações (pós-marcos) para remoção; a ANP não usa massa e normatiza cortes abaixo do leito e obrigações por faixas de profundidade, além de critérios operacionais. \\

Resíduos e limpeza do sítio &
Ambas miram minimização de impacto e remoção de fontes de risco/poluição. &
A ANP traz regra operacional objetiva de limpeza do leito (materiais/resíduos não biogênicos, distâncias/raios); a IMO trata o tema mais por critérios de avaliação (p.ex., deterioração e produtos residuais). \\

Governança documental e transparência &
A decisão deve ser suportada por documentação e justificativa. &
A ANP é muito mais prescritiva em pipeline (EJD, PDI, aprovações, relatórios, publicidade); a IMO é mais enxuta e focada no núcleo marítimo (autorização, avisos, segurança). \\

Escopo (terra, bens, contratos) &
— &
A ANP alcança também instalações terrestres e temas de reversão/alienação de bens e dinâmicas contratuais; a IMO concentra-se na remoção/gestão de estruturas offshore sob a ótica da navegação e do uso do mar. \\

Relação explícita entre as normas &
A lógica “remoção como regra + exceções justificadas” é compatível e alinhada. &
A ANP internaliza a lógica IMO no arcabouço institucional brasileiro (ANP + órgão ambiental + autoridade marítima) e acrescenta detalhamento procedimental e de enforcement. \\

\end{longtable}
\normalsize

\end{landscape}